\documentclass[11pt]{article}
% Standalone preamble for web-only supplementary notes.
\usepackage[T1]{fontenc}
\usepackage{amsmath,amssymb,amsthm}
\usepackage{hyperref}
\usepackage{cleveref}
\DeclareMathOperator{\Map}{Map}
\DeclareMathOperator{\Fun}{Fun}
\newcommand{\iso}{\xrightarrow{\sim}}
\newtheorem{lemma}{Supplementary Lemma}
\newtheorem{corollary}[lemma]{Supplementary Corollary}
\crefname{lemma}{Supplementary Lemma}{Supplementary Lemmas}
\Crefname{lemma}{Supplementary Lemma}{Supplementary Lemmas}
\crefname{corollary}{Supplementary Corollary}{Supplementary Corollaries}
\Crefname{corollary}{Supplementary Corollary}{Supplementary Corollaries}

\begin{document}

\section*{Why the endpoints matter}

The proof that \href{equivalences.html\#lem:Whiskering_By_Equivalence}{whiskering
by an equivalence induces an equivalence} uses unitors, associators, and
identifications of functors. Consequently, some intermediate comparisons have
different endpoints. The relevant elementary fact is that composition with
a fixed identification gives an equivalence between identification animae.

We work with the absolute categories and the structures of Sections 1.1--1.3.
Write $\beta\cdot\alpha$ for vertical composition, with $\alpha$ applied
first. Whiskering is denoted by $\star$. An expression such as
$\beta\cdot-$ is a functor: the fixed identification is used as a constant
input to vertical composition.

\begin{lemma}
\label[lemma]{lem:Supplementary_Endpoint_Multiplication}
For functors $f,g,h\colon C\to D$ and an identification $\beta\colon g=h$,
the functor
\[
 L_\beta\colon(f=g)\longrightarrow(f=h),\qquad
 L_\beta(\sigma)=\beta\cdot\sigma,
\]
is an equivalence, with inverse given by composition with $\beta^{-1}$.
Likewise, for $\alpha\colon f=g$, the functor
\[
 R_\alpha\colon(g=h)\longrightarrow(f=h),\qquad
 R_\alpha(\tau)=\tau\cdot\alpha,
\]
is an equivalence, with inverse given by composition with $\alpha^{-1}$.
\end{lemma}
\begin{proof}[Proof idea]
For the first assertion, the two inverse calculations are
\[
 \beta^{-1}\cdot(\beta\cdot\sigma)=\sigma,
 \qquad
 \beta\cdot(\beta^{-1}\cdot\tau)=\tau.
\]
Reassociate, apply the corresponding inverse comparison, and remove the
identity identification by a unit comparison. These calculations hold as
identifications of functors on the relevant identification animae: apply the
vertical laws to the universal input, keeping the fixed identification as a
constant family. The right-composition assertion uses the right unit and
the same inverse laws. Evaluating the composite functors also uses the
substitution comparisons from
\href{supplementary-specializing-coherence-laws.html}{Specializing coherence laws}.
\end{proof}

\begin{lemma}
\label[lemma]{lem:Supplementary_Endpoint_Square}
Let $f,g,f',g'\colon C\to D$, with identifications
$\alpha\colon f=f'$ and $\beta\colon g=g'$. Suppose that functors
\[
 u\colon X\to(f=g),\qquad v\colon X\to(f'=g')
\]
are equipped with an identification of functors on $X$
\[
 \overline\beta\cdot u=v\cdot\overline\alpha,
\]
where the bars denote constant families. Then $u$ is an equivalence if and
only if $v$ is an equivalence.
\end{lemma}
\begin{proof}[Proof idea]
After the substitution comparisons, the given identification reads
$L_\beta\circ u=R_\alpha\circ v$, with common target $(f=g')$.
Both $L_\beta$ and $R_\alpha$ are equivalences by
\Cref{lem:Supplementary_Endpoint_Multiplication}.
Apply the two-out-of-three property and invariance of equivalence under
identification of functors, in either direction.
\end{proof}

\section*{Whiskering by an equivalence}

\begin{corollary}
\label[corollary]{cor:Supplementary_Whiskering_Equivalence}
If $F\colon C\to D$ is an equivalence, then both functors
\[
 F\star-\colon(f=g)\longrightarrow(F\circ f=F\circ g)
\]
for $f,g\colon B\to C$, and
\[
 -\star F\colon(h=k)\longrightarrow(h\circ F=k\circ F)
\]
for $h,k\colon D\to E$, are equivalences.
\end{corollary}
\begin{proof}[Proof idea]
Choose an inverse $G$ of $F$. For postwhiskering, consider three successive
functors on identification animae: postwhiskering by $F$, then by $G$, then
by $F$ again. The first two compose to an equivalence: the associator
comparison relates their composite to postwhiskering by $GF$, the supplied
identification $GF=\mathrm{id}_C$ relates that to postwhiskering by the
identity, and the unitor comparison relates the latter to the identity
functor on $(f=g)$. At each step,
\Cref{lem:Supplementary_Endpoint_Square} handles the changing endpoints.
The second and third functors also compose to an equivalence, using
$FG=\mathrm{id}_D$. Two-out-of-six now implies that the first functor is
an equivalence. The prewhiskering proof uses the corresponding comparisons
on the other side.

The comparisons for composite and identity whiskering are displayed in
\href{supplementary-whiskering-and-horizontal-composition.html}{Whiskering and
horizontal composition}. Replacing $GF$ by the identity uses fixed-outer
interchange; the analogous change for prewhiskering uses fixed-inner
interchange.
\end{proof}

\section*{Lifting a specified identification}

\begin{lemma}
\label[lemma]{lem:Supplementary_Whiskering_Lift}
Let $F\colon C\to D$ be an equivalence and $f,g\colon B\to C$.
Every specified identification $\delta\colon F\circ f=F\circ g$ admits
an identification $\widetilde\delta\colon f=g$ together with an
identification
\[
 F\star\widetilde\delta=\delta.
\]
Moreover, for $\sigma,\tau\colon f=g$ and a specified identification
$p\colon F\star\sigma=F\star\tau$, there is an identification
$\widetilde p\colon\sigma=\tau$ equipped with a comparison between its
image under postwhiskering by $F\star-$ and $p$.
The corresponding assertions hold for prewhiskering by $F$.
\end{lemma}
\begin{proof}[Proof idea]
Apply a chosen inverse of the equivalence $F\star-$ to the absolute object
$\delta$. Its inverse comparison, together with the associator and unitor,
gives the specified image comparison. For the higher assertion, apply
\Cref{cor:Supplementary_Whiskering_Equivalence} to the equivalence $F\star-$
itself, and repeat the same lifting construction on the resulting
identification animae.
\end{proof}

The image comparison is part of the conclusion. It is needed, for example,
when a lifted identification must fill a prescribed square. Keeping it
avoids replacing that square by an unspecified comparison with the same
endpoints.

\end{document}
